%!TEX program = lualatex
\documentclass[a4paper,11pt]{article}

\usepackage[brazilian]{babel}
\usepackage{fontspec}
\setmainfont{Latin Modern Roman}
\usepackage[a4paper,top=1.35cm,bottom=1.35cm,left=1.65cm,right=1.65cm]{geometry}
\usepackage{graphicx}
\usepackage{tikz}
\usetikzlibrary{arrows.meta,positioning,calc,shapes.geometric}
\usepackage{array}
\usepackage{tabularx}
\usepackage{booktabs}
\usepackage{ragged2e}
\usepackage{microtype}
\pagestyle{empty}
\setlength{\parindent}{0pt}

\newcommand{\campolinha}[1]{\textbf{#1}\rule{8.6cm}{0.4pt}}

\begin{document}

\begin{minipage}[t]{0.78\textwidth}
\vspace{0pt}
\textbf{Universidade do Estado do Rio de Janeiro -- UERJ}\\
Faculdade de Geologia -- FGEL\\
Departamento de Estratigrafia e Paleontologia -- DEPA\\
Disciplina: \rule{7.2cm}{0.4pt}\\
Professora: Renata M. A. Gama\\
Exercício 1 -- Ciclo das Rochas\\[0.25em]
\textbf{ALUNO(A):} \rule{8.7cm}{0.4pt}
\end{minipage}%
\hfill
\begin{minipage}[t]{0.18\textwidth}
\vspace{-0.2cm}
\centering
\includegraphics[width=2.75cm]{uerj.png}
\end{minipage}

\vspace{0.45cm}

\begin{center}
\begin{tikzpicture}[
    x=1cm,y=1cm,
    >=Latex,
    caixa/.style={draw, minimum width=2.05cm, minimum height=0.72cm, fill=white, inner sep=2pt},
    processo/.style={draw, minimum width=2.15cm, minimum height=0.72cm, fill=white, inner sep=2pt},
    seta/.style={->, very thick},
    setaPont/.style={->, thick, dotted},
    rocha/.style={draw=none, fill=gray!28},
    sed/.style={draw=none, fill=gray!18},
    magma/.style={draw, rounded corners=8pt, fill=gray!8}
]

% --- formas geológicas simplificadas ---
% bloco metamórfico à esquerda
\path[rocha] (0.2,1.2) -- (3.4,1.2) -- (3.45,5.0) -- (2.35,5.45) -- (1.15,4.6) -- (0.2,4.25) -- cycle;
\path[fill=gray!42] (1.05,4.35) -- (2.25,5.15) -- (3.35,4.75) -- (3.25,3.35) -- (2.15,3.0) -- cycle;

% estratos sedimentares
\draw[thick] (0.0,5.55) .. controls (1.2,5.5) and (1.8,6.15) .. (3.0,6.35);
\draw[thick] (0.0,5.35) .. controls (1.2,5.3) and (1.8,5.95) .. (3.0,6.15);
\draw[thick] (0.0,5.15) .. controls (1.2,5.1) and (1.8,5.75) .. (3.0,5.95);
\draw[thick] (3.0,6.35) -- (8.7,6.35);
\draw[thick] (3.0,6.15) -- (8.7,6.15);
\draw[thick] (3.0,5.95) -- (8.7,5.95);

% corpo intrusivo
\path[fill=gray!12, draw, thick] (5.05,5.95) .. controls (5.0,6.8) and (5.4,7.25) .. (6.1,7.2)
  .. controls (6.65,7.15) and (6.95,6.65) .. (7.0,5.95) -- cycle;
\foreach \x/\y in {5.3/6.3,5.6/6.6,5.9/6.15,6.2/6.55,6.5/6.25,5.45/6.95,6.35/6.9}
  \fill (\x,\y) circle (1.1pt);

% conduto e câmara magmática
\draw[very thick, gray!55] (7.05,2.2) .. controls (7.25,3.1) and (7.2,4.3) .. (6.5,5.9);
\path[magma] (6.35,1.25) .. controls (5.7,1.25) and (5.45,1.55) .. (5.55,2.0)
  .. controls (5.7,2.5) and (6.45,2.8) .. (7.25,2.55)
  .. controls (7.9,2.35) and (8.05,1.7) .. (7.55,1.4)
  .. controls (7.2,1.18) and (6.8,1.18) .. cycle;

% corpo intrusivo profundo
\path[fill=gray!12, draw, thick] (6.0,3.25) .. controls (5.6,3.55) and (5.65,4.3) .. (6.15,4.55)
 .. controls (6.8,4.85) and (7.35,4.4) .. (7.25,3.8)
 .. controls (7.15,3.2) and (6.55,3.0) .. cycle;
\foreach \x/\y in {6.0/3.7,6.3/4.1,6.65/3.65,6.85/4.25,6.45/4.45}
  \fill (\x,\y) circle (1.1pt);

% vulcão
\path[fill=gray!38,draw,thick] (8.0,6.35) -- (8.55,7.3) -- (9.05,6.35) -- cycle;
\draw[very thick,gray!55] (8.55,2.55) -- (8.55,7.28);
\fill (8.45,7.65) circle (2pt);
\fill (8.65,7.85) circle (1.7pt);
\fill (8.82,7.55) circle (1.5pt);

% --- caixas de processos (vazias) ---
\node[processo] (p1) at (2.2,8.8) {};
\node[processo] (p2) at (4.5,8.8) {};
\node[processo] (p3) at (7.25,8.8) {};
\node[processo] (p4) at (1.2,7.75) {};
\node[processo] (p5) at (0.55,6.7) {};
\node[processo] (p6) at (0.7,4.85) {};
\node[processo] (p7) at (0.95,3.55) {};
\node[processo] (p8) at (1.7,2.3) {};
\node[processo] (p9) at (4.4,1.2) {\textbf{FUSÃO}};
\node[processo] (p10) at (6.7,2.0) {};
\node[processo] (p11) at (5.8,4.15) {};
\node[processo] (p12) at (6.4,5.3) {};
\node[processo] (p13) at (8.15,7.55) {};

% --- conexões principais ---
\draw[setaPont] (p3.west) -- (p2.east);
\draw[setaPont] (p2.west) -- (p1.east);
\draw[setaPont] (p1.west) -- ++(-0.65,0) -- (p4.north);
\draw[setaPont] (p4.south) -- (p5.north);
\draw[setaPont] (p5.south) -- (p6.north);
\draw[setaPont] (p6.south) -- (p7.north);
\draw[setaPont] (p7.south) -- (p8.north);
\draw[setaPont] (p8.east) -- (p9.west);
\draw[setaPont] (p9.east) -- (p10.west);
\draw[setaPont] (p10.north) -- (p11.south);
\draw[setaPont] (p11.north) -- (p12.south);
\draw[setaPont] (p12.north) -- (p2.south);
\draw[setaPont] (p10.east) -- ++(1.85,0) -- (p13.south);
\draw[setaPont] (p13.north) -- (p3.south);

% ligações diagonais representando rotas alternativas
\draw[setaPont] (p3.south west) -- (2.55,5.55);
\draw[setaPont] (p12.west) -- (2.85,3.75);
\draw[setaPont] (p11.west) -- (2.55,2.75);
\draw[setaPont] (p2.south west) -- (2.15,4.15);

\end{tikzpicture}
\end{center}

\vspace{0.15cm}

\textbf{Preencha as lacunas com os eventos citados abaixo, de acordo com a ordem do trajeto percorrido pela rocha.}

\vspace{0.25cm}

\renewcommand{\arraystretch}{1.15}
\begin{tabularx}{\textwidth}{|>{\RaggedRight\arraybackslash}X|>{\RaggedRight\arraybackslash}X|}
\hline
\textbf{Processos} & \textbf{Produtos} \\
\hline
Erosão / Transporte & Rochas metamórficas \\
\hline
Intemperismo & Magma \\
\hline
Deposição e diagênese & Sedimentos \\
\hline
Soerguimento e exposição & Rochas ígneas intrusivas \\
\hline
Metamorfismo & Rochas ígneas extrusivas \\
\hline
Fusão & Rochas sedimentares \\
\hline
Intrusão / Cristalização & \\
\hline
Vulcanismo / Consolidação & \\
\hline
\end{tabularx}

\end{document}
